\chapter{Iterating Inside It: Krylov Methods on the Free Block}
\label{ch:krylov}

Chapters~\ref{ch:walking} and~\ref{ch:guessing} took the inner linear solve for
granted. This chapter takes it seriously.

The observation that organises the chapter is this. Once a free set $\Free$ is
fixed, the bound-constrained problem reduces to the unconstrained SPD system
$A_\Free x_\Free = b_\Free$---and the canonical solver for an SPD system is the
conjugate gradient method of Hestenes and Stiefel \cite{hestenes1952}. CG needs
only the action $v \mapsto A_\Free v$, never the matrix. So if $A$ arrives in
factored form---$A = M\T M$ for a data matrix $M$, which is exactly what a Gram
matrix or a sample covariance is---then the whole method can run without ever
assembling an $n \times n$ object.

CG solves linear systems; it does not, unaided, respect $x \ge 0$. Wrapping it in
the active-set loop of Chapter~\ref{ch:guessing} supplies the missing layer. The
composition is not free: the loop's termination theorem assumes exact inner solves
and CG returns an approximation. Section~\ref{sec:inexact} closes that gap, and it
is the most instructive argument in the chapter.

\section{Conjugate gradients, in one page}

For SPD $A$ and right-hand side $b$, CG builds its $k$-th iterate as the minimiser
of the energy $\tfrac12 x\T A x - b\T x$ over the Krylov subspace
$\Krylov_k(A,b) = \mathrm{span}\{b, Ab, \dots, A^{k-1}b\}$, and does so with one
matrix--vector product, a handful of vector operations, and $O(n)$ storage per
iteration. That it can optimise over a growing subspace at constant cost per step
is the content of the three-term recurrence, and any of
\cite{golub2013,greenbaum1997,trefethenbau1997} derives it.
\index{conjugate gradient method}

\begin{proposition}[CG convergence]
\label{prop:cg}
Let $A_\Free$ be SPD with condition number $\kappa = \kappa(A_\Free)$, and let $x_k$
be the $k$-th CG iterate for $A_\Free x_\Free = b_\Free$. Then
\[
  \frac{\norm{e_k}_{A_\Free}}{\norm{e_0}_{A_\Free}}
  \;\le\;
  2\left(\frac{\sqrt{\kappa}-1}{\sqrt{\kappa}+1}\right)^{\!k},
\]
where $\norm{e}_{A_\Free} = (e\T A_\Free e)^{1/2}$ and
$e = x_\Free - x_k$.
\end{proposition}

Reaching a fixed relative energy-norm tolerance therefore costs $O(\sqrt\kappa)$
iterations, and clustered spectra do better still---the superlinear phase that a
Chebyshev bound cannot see. Two consequences frame the rest of the chapter.

First, by Proposition~\ref{prop:interlace} the bound applies \emph{uniformly over
every free set the loop might visit}, since $\kappa(A_\Free) \le \kappa(A)$. One
does not need to know which free sets those are.

Second, $\sqrt\kappa$ is the whole reason to keep a Krylov method rather than a
simpler iteration. Section~\ref{sec:projgrad} makes the comparison precise.

\section{The matrix-free operator}
\label{sec:operator-two}

CG needs only $v \mapsto A_\Free v$. When $A = M\T M$ this is
\[
  A_\Free v = M_\Free\T(M_\Free\,v), \qquad M_\Free = M_{:,\Free},
\]
two products at $O(m\,|\Free|)$ with \emph{no} $O(n^2)$ storage---the standard
CG-on-the-normal-equations arrangement \cite{golub2013,paige1975}, on a column
subset that changes between outer steps.
\index{matrix-free}

The whole method needs exactly two operations of $A$, and it is worth naming them
because Chapter~\ref{ch:operator} builds an abstraction on precisely this list.

\begin{enumerate}
\item \textbf{The free-block application} $v \mapsto A_\Free v$, which is all CG
      needs to solve the inner system.
\item \textbf{The cross product}
      $x_\Free \mapsto A_{\Bound,\Free}\,x_\Free$, which supplies the reduced
      gradient $s_i = (Ax)_i - b_i$ at the bound indices (where $x_\Bound = 0$),
      and hence the dual-feasibility test.
\end{enumerate}

Any implementation realising these two runs the loop, its termination proof, and its
convergence analysis unchanged. A dense $A$ slices them directly. The Gram operator
$A = M\T M$ evaluates them from $M$ without forming $A$. A diagonal-plus-low-rank
operator $A = \diag(d) + U\Delta U\T$ inverts the free block by the Woodbury
identity---a \emph{direct} inner solve that bypasses CG entirely, which is what
Chapter~\ref{ch:operator} exploits.

\section{Adding a linear equality}

Many applications carry a linear balance system alongside the bounds:
\begin{equation}
  \min_{x}\; \tfrac12 x\T A x - b\T x
  \quad\text{subject to}\quad Bx = c,\; x \ge 0,
  \qquad B \in \R^{p\times n} \text{ of full row rank}.
  \label{eq:eqp}
\end{equation}
The budget constraint $\onevec\T x = \beta$ is the case $p = 1$; group budgets,
factor neutralities, and netting rules are $p > 1$ at the same cost structure.

Naively this makes the free-set system the indefinite saddle system
of~\eqref{eq:saddle}, which CG cannot touch. The remedy is the Schur elimination of
Chapter~\ref{ch:freeblock}, and it is worth writing out because it is what keeps
every solve SPD.

Stationarity gives $x_\Free = A_\Free^{-1}(b_\Free + B_\Free\T\lambda) = v_0 + V_1\lambda$
in terms of the SPD solves
\[
  v_0 = A_\Free^{-1} b_\Free,
  \qquad
  V_1 = A_\Free^{-1} B_\Free\T \quad(\text{its $p$ columns, one CG solve each}).
\]
Imposing $B_\Free x_\Free = c$ fixes the multiplier through the $p\times p$ system
\[
  S\,\lambda = c - B_\Free v_0,
  \qquad
  S = B_\Free V_1 = B_\Free A_\Free^{-1} B_\Free\T \succ 0,
  \qquad
  x_\Free = v_0 + V_1\lambda .
\]
So each outer step costs $p+1$ matrix-free CG solves sharing one operator, plus a
$p\times p$ dense Cholesky of negligible cost. \emph{The multipliers are recovered
analytically and never appear in the linear solve}, so every inner system stays SPD
and Proposition~\ref{prop:cg} applies to each.

For $p = 1$, $B = \onevec\T$, $c = \beta$ this collapses to two solves
$v_0, v_1 = A_\Free^{-1}\onevec$ with the scalar
$\lambda = (\beta - \onevec\T v_0)/(\onevec\T v_1)$; and the pure-normalisation case
$b = 0$, $\beta = 1$ collapses further to the single solve
$x_\Free = v_1/(\onevec\T v_1)$. That last formula is the entire long-only
minimum-variance portfolio, and Chapter~\ref{ch:operator} returns to it.

\begin{remark}[What this buys over feasible-point methods]
\label{rem:feasible}
There is a well-developed family of \emph{feasible-point} bound-constrained
solvers---Bertsekas' projected Newton method \cite{bertsekas1982}, Moré and
Toraldo's GPCG \cite{moretoraldo1991}, Dostál's proportioning and MPRGP algorithms
\cite{dostal1997,dostalschoberl2005}---in which every iterate satisfies the bounds
and the working face changes through gradient projections. They are excellent
methods and their rates are stated in $\kappa$ as well.

The distinction that matters here is the equality $Bx = c$. Those methods rely on a
closed-form projection onto the feasible set, which the orthant supplies but
$\{x\ge0 : Bx=c\}$ does not; they reach a general linear equality only by wrapping
an outer augmented-Lagrangian loop \cite{dostal2006smalbe}, at the price of an extra
level of iteration and of the exact certificate. The loop of this chapter takes
$Bx = c$ directly through the Schur elimination above.

The trade runs the other way too, and honesty requires stating it: this loop treats
only the non-negative orthant $x \ge 0$, not the general two-sided box
$\ell \le x \le u$ that those methods support. A two-state complementary basis
cannot express two-sided bounds; that would need a mixed-complementarity
formulation.
\end{remark}

\section{The projected-gradient reference}
\label{sec:projgrad}

To see what the Krylov subspace buys, compare against a method that forgoes it.
Projected gradient iterates
\[
  x_{k+1} = \Pi_\Omega\!\left(x_k - \tau\,\nabla f(x_k)\right),
  \qquad \nabla f(x) = A x - b,
\]
with $\tau = 1/L$, $L = \lambda_{\max}(A)$. For the orthant
$\Pi_\Omega(y) = \max(y,0)$ componentwise; for the simplex slice
$\{x \ge 0 : \onevec\T x = \beta\}$ the projection is $O(n\log n)$ by a single
sort-and-threshold pass \cite{duchi2008,condat2016}. There is no outer loop and no
linear solve.
\index{projected gradient}

\begin{proposition}[Projected-gradient convergence]
\label{prop:proj}
Let $A \succ 0$, so $f$ is $\mu$-strongly convex and $L$-smooth with
$\mu = \lambda_{\min}(A)$, $L = \lambda_{\max}(A)$. Projected gradient with
$\tau = 1/L$ satisfies
$\norm{x_k - x^\star}^2 \le (1 - 1/\kappa(A))^{k}\norm{x_0 - x^\star}^2$.
\end{proposition}

\begin{proof}[Sketch]
$x^\star$ is a fixed point; non-expansiveness of $\Pi_\Omega$ and co-coercivity of
the $L$-smooth gradient give
$\norm{x_{k+1}-x^\star}^2 \le \norm{x_k-x^\star}^2 -
\tau(2-\tau L)\inner{\nabla f(x_k)-\nabla f(x^\star)}{x_k-x^\star}$; setting
$\tau = 1/L$ and applying strong convexity gives the claim.
\end{proof}

So the iteration count is $O(\kappa)$ against $O(\sqrt\kappa)$ for CG, at the same
per-step cost of one operator application. On ill-conditioned problems that
$\sqrt\kappa$ gap is a runtime gap. Nesterov acceleration \cite{nesterov2004,beck2009}
would lift the projected-gradient rate to $O(\sqrt\kappa)$, but CG attains that rate
\emph{optimally} within the Krylov subspace: each CG iterate minimises $f$ over
$\Krylov_k$, a strictly richer search space than two-point momentum, and typically
with a smaller constant.

The trade is otherwise in the other direction, and the honest summary is that
projected gradient is loop-free, needs no complementarity certificate, and is the
simpler method when $\kappa$ is small or moderate accuracy suffices. It also loses
its per-step simplicity the moment $p > 1$: projecting onto
$\{x\ge0 : Bx=c\}$ for general $B$ is itself a non-negative quadratic program.

\section{Inexactness: composing the two guarantees}
\label{sec:inexact}

We now close the gap flagged at the end of Chapter~\ref{ch:guessing}.

\begin{lemma}[Inexact inner solves]
\label{lem:inexact}
\index{decision margin}
Fix the bound-constrained problem~\eqref{eq:bqp}. For a free set $\Free$ let
$\hat x(\Free)$ be the exact solution of the inner system extended by zeros, and
$\hat s(\Free) = A\hat x(\Free) - b$. Let $\Free_0 = \{1,\dots,n\}, \dots, \Free_T$
be the free-set trajectory of Algorithm~\ref{alg:elim} run with \emph{exact} inner
solves, finite by Theorem~\ref{thm:termination}. Define the \emph{decision margin}
\[
  \mu = \min_{0 \le t \le T} \min\Bigl(
    \{ |\hat x_i(\Free_t)| : i \in \Free_t,\ \hat x_i \ne 0 \}
    \cup
    \{ |\hat s_i(\Free_t)| : i \notin \Free_t,\ \hat s_i \ne 0 \}
  \Bigr) > 0,
\]
the smallest nonzero magnitude any primal or dual test encounters along the exact
trajectory. Suppose $0 < \varepsilon < \mu$ and every inner CG call is stopped at a
residual $\rho = \norm{b_\Free - A_\Free\tilde x_\Free}_2$ obeying
\begin{equation}
  \rho\left(\frac{1 + \lambda_{\max}(A)}{\lambda_{\min}(A)}\right)
  \;<\; \min\bigl(\varepsilon,\ \mu - \varepsilon\bigr).
  \label{eq:inexactcond}
\end{equation}
Then at every outer step the violator sets computed from the CG iterate coincide
with those computed from the exact solve on the same free set. The inexact loop
therefore traverses exactly the trajectory $\Free_0,\dots,\Free_T$,
Theorem~\ref{thm:termination} applies verbatim, and
$\norm{\tilde x - x^\star}_\infty \le \rho/\lambda_{\min}(A)$.
\end{lemma}

\begin{proof}
Write $r = b_\Free - A_\Free\tilde x_\Free$, so
$\tilde x_\Free - \hat x_\Free = -A_\Free^{-1}r$. By Cauchy interlacing
(Proposition~\ref{prop:interlace}), $\lambda_{\min}(A_\Free) \ge \lambda_{\min}(A)$,
hence
$\delta_x := \norm{\tilde x_\Free - \hat x_\Free}_2 \le \rho/\lambda_{\min}(A)$,
which also bounds each coordinate. For $i \notin \Free$,
$\tilde s_i - \hat s_i = A_{i,\Free}(\tilde x_\Free - \hat x_\Free)$, so by
Cauchy--Schwarz $\delta_s \le \lambda_{\max}(A)\delta_x$. Both are bounded by
$\delta := \rho(1+\lambda_{\max}(A))/\lambda_{\min}(A) < \min(\varepsilon,\mu-\varepsilon)$.

Fix $i \in \Free_t$. If $\hat x_i \ge 0$ then $\hat x_i = 0$ or $\hat x_i \ge \mu$;
either way $\tilde x_i \ge -\delta > -\varepsilon$ and the inexact test does not flag
$i$, matching the exact test. If $\hat x_i < 0$ then $\hat x_i \le -\mu$, so
$\tilde x_i \le -\mu + \delta < -\varepsilon$ and both tests flag $i$. The same
argument with $\hat s_i, \tilde s_i, \delta_s$ covers the dual test at every bound
index. Hence the violator sets agree at step $t$; since both loops start from the
same $\Free_0$ and the counters are deterministic functions of the violator
sequence, induction on $t$ gives the same trajectory and the same stopping step. At
termination the exact iterate is $x^\star$, so
$\norm{\tilde x - x^\star}_2 = \delta_x \le \rho/\lambda_{\min}(A)$.
\end{proof}

Read the structure of that argument, because it generalises. One does not prove a
new termination theorem for the approximate method. One proves that the approximate
method makes the \emph{same decisions} as the exact method, and then inherits the
exact method's theorem unchanged. Whenever an algorithm's control flow depends on
finitely many sign tests, and those tests have a margin, this is the argument
available.

The corresponding statement for the equality-augmented problem~\eqref{eq:eqp} needs
one extra step, because the multiplier $\lambda$ is itself recovered from $p+1$
inexact solves through the Schur complement rather than read off a single residual.
Propagating the perturbation through $\tilde S^{-1}$ by a Neumann series gives a
bound of the same shape, and the sign-test argument then runs verbatim.

\begin{remark}[The guarantee attaches to the outer loop]
Nothing in Theorem~\ref{thm:termination} mentions CG. A \emph{direct} factorisation
may replace the Krylov inner solve unchanged, and then no inexactness condition
needs verifying at all. Chapter~\ref{ch:operator} does exactly this with a Woodbury
solve. The separation is worth naming: the outer loop owns the combinatorics and the
certificate; the inner solver owns the linear algebra; and Lemma~\ref{lem:inexact} is
the interface contract between them.
\end{remark}

\section{Warm starts}

For a parametric family---a sequence of right-hand sides $b(\theta_1), b(\theta_2),
\dots$, or a swept normalisation $\beta$---consecutive minimisers usually differ in
only a few active constraints. Passing the previous problem's $(\Free, x)$ as the
initial state reduces both counts: the outer count, because the free set needs fewer
adjustments, and the inner count, because CG starts from a small initial residual.

\begin{proposition}[Warm starts across a support-stable step]
\label{prop:warm}
Let $b' = b + \delta$ and suppose the free set $\Free$ returned for $b$ is also
optimal for $b'$. Then Algorithm~\ref{alg:elim} started from $\Free$ terminates in a
\emph{single} outer step, and its one CG solve, warm-started at the previous
solution, reaches energy-norm accuracy $\tau$ in at most
\[
  k \;\le\; \left\lceil \frac{\sqrt{\kappa}+1}{2}\,
  \ln\!\frac{2\,\norm{\delta}_2}{\tau\sqrt{\lambda_{\min}(A)}} \right\rceil
  \text{ iterations.}
\]
\end{proposition}

\begin{proof}
The single solve returns $\hat x_\Free \ge 0$ with $s \ge 0$ on the bound set, so
both violator sets are empty and the loop exits, certified by Step~2 of
Theorem~\ref{thm:termination}'s proof. The warm start's initial error is
$e_0 = A_\Free^{-1}\delta_\Free$, whose energy norm is
$(\delta_\Free\T A_\Free^{-1}\delta_\Free)^{1/2} \le \norm{\delta}_2/\sqrt{\lambda_{\min}(A)}$
by interlacing. Proposition~\ref{prop:cg} with
$\rho_{\mathrm{cg}} = (\sqrt\kappa-1)/(\sqrt\kappa+1)$ and
$\ln(1/\rho_{\mathrm{cg}}) \ge 2/(\sqrt\kappa+1)$ gives the bound.
\end{proof}

The interesting feature is that the inner count scales with the logarithm of the
\emph{parameter step} $\norm{\delta}$, not of the solution norm: a finer grid is
cheaper per point. When the support drifts, no comparable a-priori bound is
available---the loop is not monotone in the free set---but empirically the
warm-started outer count tracks the support drift $|\Free \triangle \Free'|$ and
collapses to one whenever the drift is zero, exactly as the proposition prescribes.

Note also which solver profits from what. A \emph{direct} inner solver profits only
from the warm free set, having no analogue of an initial guess. A matrix-free CG
inner solver profits from both. That is where warm-started parametric sweeps obtain
their largest speed-ups, and Chapter~\ref{ch:warmstart} develops the theme.

\begin{notes}
The construction of this chapter is \cite{schmelzer2026nncg}, released as the
open-source \texttt{nncg} package; Lemma~\ref{lem:inexact} and its
equality-augmented corollary are its contribution, the exact-solve termination count
being Júdice and Pires' \cite{judice1994}. CG is \cite{hestenes1952};
\cite{greenbaum1997} and \cite{trefethenbau1997} are the textbook treatments of its
convergence, and \cite{golub2013} of the normal-equations arrangement. On the NNLS
side, Bro and De Jong's fast NNLS \cite{brodejong1997} accelerates Lawson--Hanson
\cite{lawson1974} by caching normal-equations cross-products, but still assembles
$M\T M$; the method here forms neither.
\end{notes}

\begin{exercises}

\begin{exercise}
Show that $A = M\T M$ is positive definite iff $M$ has full column rank, and that
the free-block operator $v \mapsto M_\Free\T(M_\Free v)$ is SPD for every $\Free$
when it is.
\end{exercise}

\begin{exercise}
Derive the $p = 1$ specialisation of the Schur elimination, and verify that
$b = 0$, $\beta = 1$ gives $x_\Free = v_1/(\onevec\T v_1)$.
\end{exercise}

\begin{exercise}
Show $S = B_\Free A_\Free^{-1} B_\Free\T \succ 0$ whenever $A_\Free \succ 0$ and
$B_\Free$ has full row rank. For $p = 1$, when can $B_\Free$ fail to have full row
rank?
\end{exercise}

\begin{exercise}
Prove that CG applied to $A_\Free x = b_\Free$ from $x_0 = 0$ produces iterates in
$\Krylov_k(A_\Free, b_\Free)$, and deduce that the exact solution is reached in at
most $|\Free|$ iterations in exact arithmetic. Why is this fact of limited practical
use?
\end{exercise}

\begin{exercise}
The decision margin $\mu$ of Lemma~\ref{lem:inexact} is defined along the exact
trajectory, which one does not know in advance. Explain why the lemma is
nevertheless useful, and propose a practical stopping rule that respects its spirit.
\end{exercise}

\begin{exercise}
Compare the total work of CG and projected gradient to reach relative accuracy
$10^{-8}$ on a problem with $\kappa = 10^{4}$, assuming one operator application per
step in both. Now repeat for $\kappa = 10^2$ and comment.
\end{exercise}

\begin{exercise}[Implementation]
Implement Algorithm~\ref{alg:elim} with a matrix-free CG inner solver on
$A = M\T M$. On a synthetic NNLS problem with a planted non-negative solution, plot
total CG iterations against the inner residual tolerance and locate the point at
which the outer trajectory changes---i.e.\ where~\eqref{eq:inexactcond} fails.
\end{exercise}

\begin{exercise}[Implementation]
Sweep the normalisation $\beta$ over a grid of $50$ values with and without warm
starting. Plot outer steps and total inner iterations for both, and check
Proposition~\ref{prop:warm}'s prediction that the inner count scales with
$\log\norm{\delta}$ by refining the grid.
\end{exercise}

\end{exercises}
